\section{Simulation Agents：从任务求解到社会过程建模}

\subsection{为什么要做 simulation}
讲者随后转向 simulation agents：不仅让 agent 解决任务，也让它在社会语境中持续互动，用于研究群体行为、信息传播与协作机制。

\begin{importantbox}{Simulation 的价值不在“像不像人”，而在“能否形成可检验假设”}
当 agent 作为可编程个体参与多主体系统时，我们可以系统地改变记忆策略、沟通规则、激励函数，观察宏观行为如何变化。
\end{importantbox}

\subsection{生成式社会的三层机制}
\begin{knowledgebox}{微观到宏观的桥梁}
\begin{enumerate}
    \item 记忆层：个体如何保留和召回经历；
    \item 决策层：个体如何在上下文中行动；
    \item 交互层：个体之间如何影响彼此并形成群体动态。
\end{enumerate}
这三层联动才构成“可研究的社会模拟系统”。
\end{knowledgebox}

\begin{warningbox}{仿真结果的外部有效性风险}
仿真系统中的规则往往由研究者设定，因此结果可能更反映“规则选择”而非真实社会机制。  
如果没有充分的现实校验，simulation 结论应被视作假设生成，而非事实证明。
\end{warningbox}

\subsection{本章小结}
Simulation agents 把 Agent 研究扩展到“系统行为层”。它们最有价值的产出不是故事，而是可以被实验检验和反驳的机制假设。


